\section{Evaluierung und Ergebnisse}
\label{sec:evaluierung}
Nach der erfolgreichen Implementierung der Klassifikationspipeline auf dem eingebetteten System wurde dieses im praktischen Einsatz evaluiert. Die nachfolgenden Abschnitte betrachten das Systemverhalten auf Basis empirischer Daten und analysieren die Interaktionslatenz, die biomechanischen Besonderheiten der Gesten-Erkennung sowie die Robustheit des Modells gegenüber veränderten Sensorpositionen.

    \subsection{Systemlatenz und Echtzeit-Verhalten}
    Eine wichtige Metrik für die Alltagstauglichkeit einer myoelektrischen Prothese ist neben der reinen Erkennungsgenauigkeit das subjektive Empfinden der Verzögerung zwischen der mentalen Intention des Nutzers und der physischen Umsetzung durch die Aktuatoren. Liegt diese Latenz hoch, wirkt das System träge und lässt sich nur schwer präzise steuern. Bei einer optimalen Positionierung der Elektroden-Armbänder weist das Gesamtsystem eine empirisch ermittelte Verzögerung von etwa 0,5 Sekunden auf.

    Diese Gesamtlatenz setzt sich aus drei sequentiellen Phasen zusammen. Da die mathematischen Features über einen gleitenden Puffer von zwei Sekunden berechnet werden, benötigt das System eine kurze Zeitspanne, bis sich die veränderten Signalamplituden im \acf{MAV} und \acf{RMS} signifikant niederschlagen. Durch die hohe Inferenzfrequenz von 10\,Hz wird die daraufhin getroffene Gestenvorhersage jedoch nahezu verzögerungsfrei in den internen Systemzustand übersetzt. Den letzten Anteil der Verzögerung bildet schließlich die mechanische Trägheit der Servomotoren, welche eine kurze physische Anlaufzeit benötigen, um die Finger der Handprothese in die Zielposition zu bringen. Mit einer Gesamtlatenz von 500\,ms liegt das System in einem Bereich der als angemessen flüssig wahrgenommen wird, ohne dass das Motorschutzfenster von einer Sekunde den Bewegungsfluss blockiert.

    \subsection{Biomechanische Signalcharakteristika und Besonderheiten}
    Bei genauerer Betrachtung der Klassifikationsergebnisse zeigt sich, dass die \textit{Schere, Stein, Papier} Gesten aus biomechanischer Sicht ideal gewählt sind. Die drei ausgewählten Gesten nutzen anatomisch komplementäre Muskelgruppen des Unterarms, was sich in markanten, leicht trennbaren Mustern über die acht Kanäle des \textsf{ADS1298} widerspiegelt. Die Geste \textit{Stein} erfordert eine volle Flexion aller Finger. Dies führt zu einer massiven, homogenen Aktivierung primär auf jenen Elektroden, die an der Innenseite des Unterarms anliegen. Aufgrund dieser eindeutigen anatomischen Charakteristik wird die Faust vom Random Forest nahezu fehlerfrei und mit maximaler Verlässlichkeit erkannt.

    Im Gegensatz dazu aktiviert die Geste \textit{Papier} primär Extensoren an der Außenseite des Unterarms, um die Finger vollständig zu strecken. Die Geste \textit{Schere} wiederum bildet eine Mischform, bei der Zeige- und Mittelfinger gestreckt, Ring- und kleiner Finger sowie der Daumen jedoch angewinkelt werden. Hierbei werden Muskelstränge der Innen- und Außenseite simultan aktiviert. Aufgrund dieser anatomischen Überschneidungen zeigt die empirische Evaluation, dass Konfusionen des Modells am ehesten zwischen den Klassen \textit{Papier} und \textit{Schere} auftreten. Wenn die Streckbewegung der Schere sehr dominant ausgeführt wird, neigt der Klassifikator dazu, die Geste fälschlicherweise als offene Hand zu interpretieren. Ein weiteres Phänomen betrifft den Übergang beim Lösen der Geste \textit{Stein} in den Ruhezustand. In dem Moment, in dem die Anspannung der inneren Beugemuskeln nachlässt, wird für einen Sekundenbruchteil oft die Geste \textit{Schere} erkannt. Diese Fehlklassifikation resultiert vermutlich aus der ungleichmäßigen Muskelentspannung und dem damit verbundenen Einschwingverhalten des Signals im abklingenden Analysefenster. Dennoch bleibt die Gesamtperformance des Systems auf einem hohen und stabilen Niveau.
    
    \subsection{Robustheit gegenüber Sensor-Repositionierung (Spatial Shift)}
    Ein kritisches Phänomen bei \ac{sEMG}-Systemen ist die Empfindlichkeit gegenüber dem erneuten Anlegen der Sensorik nach einer Nutzungspause. In der alltäglichen Praxis ist es für einen Nutzer kaum verlässlich möglich, die Elektroden-Armbänder nach dem Ablegen wieder exakt in derselben Position und Rotationsausrichtung auf dem Unterarm zu platzieren. Die Evaluierung dieses Szenarios ergab, dass die Genauigkeit des Modells von den ursprünglichen 90\,\% auf Werte zwischen 60\,\% und 70\,\% abfällt, wenn das Armband rein nach Augenmaß und ohne technische Hilfsmittel neu positioniert wird. In diesem Fall passen die gelernten Verzweigungsschwellenwerte der Entscheidungsbäume nicht mehr perfekt zu den nun leicht verschobenen Kanälen. Gelingt eine exakte Repositionierung, bleibt die hohe Ausgangsgenauigkeit erhalten, was die Validität des Modells unterstreicht. Ein Referenzfoto hilft beim Wiederanlegen, ist jedoch nicht ausreichend um eine exakte Positionierung verlässlich zu ermöglichen.

    Da eine perfekte Platzierung im Alltag unrealistisch ist, wurde untersucht wie die Genauigkeit des Systems wiederhergestellt werden kann. Es zeigte sich, dass bereits ein minimales Nachtraining (Fine-tuning) mit nur etwa 30 neuen Datenpunkten in der veränderten Position ausreicht um die Klassifikationsgenauigkeit des Modells direkt wieder auf das hohe Niveau von 80\,\% bis 90\,\% anzuheben. Dies entspricht einem kleinen Datensatz von zehn Gesten pro Klasse, welcher mit dem in \autoref{subsec:data_acquisition} beschriebenen Verfahren innerhalb von 90 Sekunden (1,5 Minuten) aufgenommen werden kann. Da auch das Modelltraining, der Modell-Export und die Programmierung des \textsf{Teensy~4.1} sehr zeiteffizient möglich ist, kann das ganze System von einem geübten Anwender innerhalb von etwa 3 Minuten auf die neue Armbandpositionen angepasst werden. Diese Beobachtung beweist, dass das trainierte Modell eine hohe Generalisierbarkeit gegenüber leichter Verschiebungen der Sensoren besitzt, die durch ein effizientes Fine-tuning kompensiert werden kann. Der mobile Einsatz der Prothese ist dadurch jedoch signifikant eingeschränkt.
